\documentclass[12pt]{article}
\usepackage[margin=2.5cm]{geometry}
\usepackage{parskip}
\usepackage{hyperref}

\title{Statement of Differences Between the Conference Version\\
and the Extended Version}
\author{Mona Rajhans \and Vishal Khawarey}
\date{}

\begin{document}
\maketitle

\noindent
The conference version of this paper~\cite{khawarey2026empirical} presented an
empirical study of adversarial robustness and SHAP-based explainability drift for
a Multi-Layer Perceptron~(MLP) classifier evaluated on two cybersecurity datasets:
a phishing URL detection dataset and the UNSW-NB15 network intrusion dataset.
Adversarial robustness was measured under the Fast Gradient Sign Method~(FGSM) and
Projected Gradient Descent~(PGD), and the Robustness Index~(RI) was introduced as
a scalar summary of the accuracy--perturbation curve.  The paper further
characterised the coupling between gradient-based feature sensitivity and SHAP
attribution drift, and demonstrated that adversarial training partially recovers
robustness on both datasets.

This extended version makes the following substantive contributions that together
constitute well over 30\% new material relative to the conference paper.

\textbf{(i) New classifier families.}
Random Forest and XGBoost are added alongside the MLP, representing the
tree-ensemble family that dominates production cybersecurity deployments but was
absent from the conference study.  Both models are evaluated under identical
experimental conditions across all datasets and attack methods.

\textbf{(ii) New attack methods.}
Three black-box attacks are introduced and evaluated: ZOO~\cite{chen2017zoo}
(finite-difference gradient approximation), Square
Attack~\cite{andriushchenko2020square} (score-based random search), and
HopSkipJump~\cite{chen2020hopskipjump} (decision-boundary estimation).  This
extends the attack surface from gradient-based white-box attacks, which are
inapplicable to non-differentiable tree models, to a fully model-agnostic
black-box evaluation framework.

\textbf{(iii) New metric: Explainability Stability Index~(ESI).}
ESI is defined as the complement of the normalised area under the SHAP-drift curve,
providing an explanation-domain analogue of RI that operates on the same
$[0,1]$ scale.  ESI is computed via TreeSHAP~\cite{lundberg2020treeshap} for
Random Forest and XGBoost, enabling direct cross-model comparison of prediction
robustness and explanation stability within a single unified framework.

\textbf{(iv) New datasets.}
The NF-ToN-IoT dataset~\cite{sarhan2021nftoniot} (1.16\,M NetFlow~v2 records,
10 features) is added as a third evaluation domain covering IoT network intrusion
detection.  HIKARI-2021~\cite{ferrag2022hikari} (555{,}278 flows, 83 features)
is added as a fourth domain covering encrypted TLS/HTTPS traffic detection,
providing additional cross-domain validation of all findings across structurally
different feature spaces.

\textbf{(v) New empirical findings.}
Evaluation reveals that ZOO produces near-degenerate gradient estimates against
XGBoost (RI\,$\approx$\,0.98) due to its piecewise-constant prediction surface,
whereas Square Attack reduces XGBoost accuracy to near-random (RI\,=\,0.36 on
Phishing).  This 0.62 gap exposes a methodological risk: attack-method selection
is not model-agnostic.  A query-budget ablation (varying both the fraction of
features sampled per gradient call from 12\% to 100\%, and the number of
iterative refinement steps from 20 to 200) confirms that XGBoost ZOO degeneracy
is budget-invariant on Phishing (RI\,=\,0.970 at all budgets, $\Delta=0.000$),
establishing it as an intrinsic property of the piecewise-constant surface
rather than an artefact of the query budget used in the main experiments.  Additionally, we demonstrate that XGBoost achieves high
prediction robustness under ZOO (RI\,$\approx$\,0.98) yet exhibits low explanation
stability (ESI\,$\approx$\,0.06--0.16 across all four datasets), showing that prediction
robustness and explanation stability are \emph{decoupled}---a high RI does not
guarantee trustworthy explanations under attack.  To confirm the RF\,>\,XGB ESI ordering is not itself a ZOO artefact, we recompute
ESI under Square Attack across all four datasets.  The ordering holds on Phishing
(RF\,=\,0.230, XGB\,=\,0.117) and UNSW-NB15 (RF\,=\,0.184, XGB\,=\,0.155), is
tied on NF-ToN-IoT, and reverses on HIKARI-2021 due to a normalisation-scale
artefact (XGBoost's absolute SHAP drift is $\approx$10$\times$ larger than RF's,
but grows more gradually, which the normalised-AUC formula scores as higher
stability).  Under Square Attack the RF\,>\,XGB ordering is confirmed on two of four datasets
(Phishing, UNSW-NB15), tied on NF-ToN-IoT, and reverses on HIKARI-2021 due to a
normalisation-scale artefact; the ordering holds on all four datasets under ZOO
(the attack used to compute the ESI values in the main results table).

\textbf{(vi) PGD step-size ablation.}
A controlled ablation over PGD step sizes $\alpha \in \{0.005, 0.01, 0.02, 0.05,
0.10\}$ confirms that the PGD\,$<$\,FGSM anomaly observed in the conference paper
is a hyperparameter sensitivity artifact on z-score normalised tabular data, not a
structural property of the model.  PGD with $\alpha=0.01$ (the conventional choice)
yields RI\,=\,0.82 versus FGSM's RI\,=\,0.68 on Phishing; the gap closes only at
$\alpha \geq 0.05$.  This provides a rigorous empirical response to reviewer
feedback on the conference version.

\textbf{(vii) Master comparative analysis.}
A comprehensive results table spanning all four datasets, three model families,
and five attack methods is introduced, enabling cross-domain and cross-model
comparison in a single view.  Twelve new figures accompany the extended experiments:
the ZOO-degeneracy accuracy curve, full attack comparison curves (ZOO, Square
Attack, HopSkipJump), the master RI/ESI heatmap (Phishing, UNSW-NB15, NF-ToN-IoT),
per-feature TreeSHAP drift heatmaps for both tree models on Phishing, NF-ToN-IoT
robustness curves, HIKARI-2021 robustness curves, the PGD step-size ablation, ESI
drift curves comparing ZOO versus Square Attack on Phishing, ESI drift curves on
UNSW-NB15 confirming the RF\,>\,XGB ordering under Square Attack,
the Square Attack adversarial training before/after curves for tree ensembles,
and the ZOO query-budget ablation (coordinate-sampling and iterative-step variants).

\textbf{(viii) Feature-level TreeSHAP drift analysis.}
A new subsection decomposes the aggregate ESI score into per-feature attribution
drift, revealing which feature types drive explanation instability for each model
family.  The analysis confirms that XGBoost's boosted-additive structure produces
higher per-feature drift than Random Forest's ensemble averaging, consistent with
the consistently lower XGBoost ESI observed across all four datasets.

\textbf{(ix) Attack-method selection guidance.}
A new discussion section introduces a two-axis framework~(gradient dependence,
query efficiency) that explains the observed attack ranking across all four
datasets and yields concrete practical recommendations: Square Attack should be
the primary benchmark for tree ensembles; ZOO alone overestimates robustness by
up to 0.62 RI; adaptive PGD step sizes are recommended for z-score normalised
tabular data.

\textbf{(x) Variance and significance study.}
A five-seed variance study (varying train/test split, model initialisation, and
evaluation subsets) is added to the Limitations section.  The ZOO degeneracy gap
is highly reproducible (mean 0.534--0.575, std\,$\leq$\,0.035 across all four
datasets), confirming the primary finding is not a lucky point estimate.  ESI
variance is higher (RF ESI std up to 0.101); the RF\,$>$\,XGB ordering holds in
all five seeds on UNSW-NB15 and in 4/5 seeds on Phishing and HIKARI-2021, but is
within noise on NF-ToN-IoT ($0.025 \pm 0.082$).

\textbf{(xi) Square Attack adversarial training for tree ensembles.}
As a black-box analogue of MLP FGSM adversarial training, Square Attack data
augmentation ($25\%$ adversarial ratio at $\varepsilon_{\max}=0.3$) is evaluated
on Phishing and UNSW-NB15.  The largest beneficiary is XGBoost on Phishing,
where Square RI rises from $0.372$ to $0.755$ ($+0.383$), nearly closing the gap
with baseline RF.  RF on UNSW-NB15 improves by $+0.156$, preventing the complete
accuracy collapse ($0.000$) observed above $\varepsilon=0.167$ in the baseline.
XGBoost on UNSW-NB15 shows cross-attack transfer ($+0.045$ ZOO, $+0.081$ HSJ),
while RF on Phishing is slightly degraded ($-0.049$), revealing that augmentation
at $\varepsilon_{\max}$ can backfire for already-robust models.  Clean accuracy
cost is at most $0.7\%$ across all conditions.

\textbf{(xii) Limitations section.}
A dedicated Limitations section discusses six scope constraints absent from the
conference version: feature-validity constraints on $\ell_\infty$ perturbations
in structured security data, dataset scope, CPU-only computational budget,
the approximation of ESI via ZOO-perturbed inputs, the variance study
results quantifying which conclusions are statistically robust, and the
scope of tree ensemble defences~(Square Attack augmentation evaluated here;
other defence strategies remain future work).

All tables, figures, and textual descriptions in this extended version are either
entirely new or substantially revised.  Results from the conference paper are
summarised in prose with explicit citation~\cite{khawarey2026empirical} rather
than reproduced verbatim.  The title, abstract, and conclusions have been rewritten
to reflect the extended scope of the work.  The Related Work section has been
expanded to position ESI explicitly against existing explanation-stability metrics
(OpenXAI RIS/ROS, Alvarez-Melis and Jaakkola sensitivity measures) and to address
the gap in tree-ensemble robustness literature.  Experiment code and result tables
are publicly available at~\cite{khawarey2026code}.

\bibliographystyle{splncs04}
\bibliography{refs}
\end{document}
